% c09 —— §14 刚性 Calabi–Yau 流形的镜像，p036
\section{刚性 Calabi–Yau 流形的镜像}

设 $\PP^8$ 中的三次超曲面为 $Z$。则 $Z$ 的 Hodge 数由 Lefschetz 定理，在中间维数（即 $p+q=7$）以外由 $h^{p,q}(Z)=\delta_{pq}$ 给出。而在中间维数，用 Griffiths 方法 [Gri69] 可知 Hodge 数为
\begin{equation}
    (h^{7,0},h^{6,1},\dots,h^{0,7})=(0,0,1,84,84,1,0,0) \tag{14.1}
\end{equation}
并且对 $H^{5,2}(Z)$ 的非零元 $\Omega$，映射
\begin{equation}
    H^1(\Theta)\to H^{4,3}(X),\qquad X\mapsto \iota_X\Omega \tag{14.2}
\end{equation}
是同构。请注意 (14.2) 与 (9.4) 的相似性。这样，具有与 Calabi–Yau 流形类似性质的 Fano 流形，在 [IM15] 中被称为 Calabi–Yau 型 Fano 流形，那里给出了暂定的定义与若干例子。容易看出 $(\cO_Z,\cO_Z(1),\cO_Z(2),\cO_Z(3),\cO_Z(4),\cO_Z(5))$ 构成 $Z$ 上的例外序列，而它的左正交 $A_Z$ 成为 3 维的 Calabi–Yau 范畴，这作为 [Kuz04, Corollary 4.3] 的特殊情形可以立即得到。据此，把虚拟的（或者说「非交换的」）Calabi–Yau 流形 $\check Y$ 定义为
\begin{equation}
    D^b\Coh\check Y := A_Z \tag{14.3}
\end{equation}
一般地，代数簇的 Hochschild 同调通过 Hochschild–Kostant–Rosenberg 同构与 Hodge 上同调以
\begin{equation}
    \mathrm{HH}_i(Z)\cong \bigoplus_{p-q=i} H^q(\Omega^p_Z) \tag{14.4}
\end{equation}
相联系，而对 $\PP^8$ 的三次超曲面 $Z$ 有
\begin{equation}
    hh_{-3}=h^{2,5}=1,\quad hh_{-2}=0,\quad hh_{-1}=h^{3,4}=84,\quad hh_0=\sum_{p=0}^{7}h^{p,p}=8, \tag{14.5}
\end{equation}
成立。对预三角微分分次范畴的半正交分解
\begin{equation}
    \cD=\langle \cD_1,\dots,\cD_n\rangle \tag{14.6}
\end{equation}
Hochschild 同调有直和分解
\begin{equation}
    \mathrm{HH}_\bullet(\cD)\cong \bigoplus_{i=1}^{n}\mathrm{HH}_\bullet(\cD) \tag{14.7}
\end{equation}
% [Tutor Note: 原文 (14.7) 右端写作 \bigoplus_i HH_\bullet(D)，按上下文（半正交分解 D=⟨D_1,…,D_n⟩ 的 Hochschild 同调分解）应为 \bigoplus_i HH_\bullet(D_i)；此处照原文保留，请审校确认。]
这一点由 [Kuz] 所知，因此
\begin{equation}
    hh_{-3}(A_Z)=hh_3(A_Z)=1,\quad hh_{-1}(A_Z)=hh_1(A_Z)=84,\quad hh_0(A_Z)=2 \tag{14.8}
\end{equation}
成立。由此可知，$\check Y$ 的（虚拟的）Hodge 数由
\begin{equation}
    \begin{pmatrix}
      & & & 1 & & & \\
      & & 0 & & 0 & & \\
      & 0 & & 0 & & 0 & \\
      1 & & 84 & & 84 & & 1 \\
      & 0 & & 0 & & 0 & \\
      & & 0 & & 0 & & \\
      & & & 1 & & &
    \end{pmatrix} \tag{14.9}
\end{equation}
给出，(9.17) 中的 $Y$ 与 $\check Y$ 成为拓扑镜像。关于刚性 Calabi–Yau 流形的镜像，也请参照 [CDP93, Sch96, BB97]。特别地，[CDP93] 不仅讨论了拓扑镜像对称，也讨论了经典镜像对称。

% CHUNK c09 DONE pages p036
